\subsection{Limitations}
\label{sec:limitations}

The sizing model governs internal element geometry. Several categories of UI concern fall outside its scope:

\begin{itemize}
  \item \textbf{Layout regions.}
  Structural containers such as dialogs, drawers, and menus have content-determined height. The model assigns them wrapping level $w = 3$ for padding and radius computation, but does not constrain their overall height. This is intentional: layout regions frame arbitrary nested content whose extent cannot be predicted by a closed-form expression.

  \item \textbf{Behavior patches.}
  Patches that apply logic without contributing visual geometry (e.g., form validation, field binding, transition orchestration) are outside the spatial model. They occupy no height and require no sizing parameters.

  \item \textbf{Separator.}
  The horizontal rule uses a fixed $1\,\text{px}$ height. This is a degenerate case that does not participate in the $U$-based model because a separator's role is to mark a boundary, not to contain content.

  \item \textbf{Custom artistic elements.}
  The model provides a baseline geometry. Elements with deliberately non-standard proportions (e.g., hero banners, decorative illustrations) may override the formula. The model does not prevent overrides; it establishes the default from which overrides are measured.

  \item \textbf{CSS border interaction.}
  The height formula assumes the CSS box model without borders. A $1\,\text{px}$ border adds $2\,\text{px}$ to rendered height ($5.56\%$ deviation at $w = 1$, $d = 1.5$). The model recommends \texttt{outline} or \texttt{box-shadow} for visual edges, but does not enforce this~\cite{css-box}.
\end{itemize}

These boundaries are explicit. The coverage proof (Table~\ref{tab:coverage-proof}) accounts for all five categories: layout regions and behavior patches are classified as n/a, the separator is classified as a fixed-height exception, and all remaining patches (54 of 69) are fully governed by the parametric formula with zero deviations.
